\chapter{First When There's Nothing --- But a Slow Glowing Dream}
\label{chap:overture}

\noindent\emph{This chapter tells the whole story in plain words---no jargon, and
no mathematics beyond a few enormous numbers. Everything in it is made precise,
named, and defended in the chapters that follow. If you read only one chapter,
read this one, and trust that each claim has its arithmetic later.}

\section*{The void}

Start with nothing. Not empty space with a clock ticking in it---really nothing:
no stars, no matter, no before or after, not even a direction to time. But
``nothing,'' in physics, is never quite still. The vacuum we have here, in our own
laboratories, is restless: strip away every particle and what remains still
shimmers with possibility. Grant the void that same restlessness---no more than
the vacuum already does---and the story tells itself.

A warning about words, though, because it matters here more than anywhere. I am
going to say ``once in a while,'' ``then,'' and ``begins,'' because our minds and
our language are built for time and cannot work without it. The void has none.
Nothing \emph{happens} in it in sequence; there is no clock, no first and next.
Picture it instead as one single, still, everlasting state whose very nature
already \emph{contains} these rare worlds---the way a held chord already contains
its overtones, sounding all at once with nothing happening in time. Worlds are not
events the void goes through; they are the faint notes inside its one unending
tone. ``Before'' and ``after'' switch on only inside a world, once its dawn has
lit an arrow. So read what follows as a translation: a timeless truth, told in the
only tense we have.

Most of those faint notes are nothing at all. But a rare few are loud enough to
cross a threshold---enough to fold up, under their own gravity, into a black hole.
And here is the first surprise. The seed you need is small---about the mass of a
mountain---and a black hole is the \emph{easiest} thing of its size to be, because
it is the most generic: the least special, most disordered object there is. The
void need not contain a finished, finely arranged universe. It need only contain a
cheap little seed, and gravity does the rest.

How rare is even that seed? About one in $10^{10^{84}}$---a one followed by a
string of zeros that itself has eighty-five digits. Vanishing in any finite patch.
But the void is endless, and even the vanishingly rare is present, somewhere, in
an endless expanse. The point was never that it is likely. The point is that it is
the \emph{cheap} way in. The eternal dawn.

\section*{The bounce, and the first dawn}

Follow the matter down into the black hole. There is the famous surface---the
horizon---past which light cannot climb back out, and where, to anyone watching
from outside, time appears to freeze. But that freezing is only the outsider's
illusion. The matter itself feels nothing special at the horizon; its own clock
keeps ticking, smoothly, all the way down. And for this first world there was no
earlier clock to borrow: time is being switched on here. What I am calling its
clock is the newborn world's own---lit at the bounce, and reaching back to that
instant as its beginning.

As it falls, it crowds together, denser and denser. Then something pushes back.
Einstein's gravity, extended by the geometer \'Elie Cartan to include the
\emph{spin} of matter, says that when matter is squeezed hard enough, that spin
resists---and at a stupendous density this resistance overwhelms gravity itself.
The collapse does not end in a point of infinite density. It \emph{bounces}.

That bounce is the dawn. It is the lowest, quietest, most ordered moment of a
brand-new universe---its beginning, its ``big bang,'' seen from the inside. There
is no waiting room of frozen time between the horizon and the bounce; the clock
runs the whole way, reaches its stillest moment at the bounce, and that moment is
time zero for the new world. The direction of time is set right there---it points
away from the bounce, toward the future, because that is the way toward more
disorder. And the spin of the original seed makes one last choice at that instant:
whether the new universe is built of matter or of antimatter. Spin too little, and
the whole thing fizzles back into the void. Spin enough, and you have made a
universe out of nothing.

\section*{The first universe is not ours}

This new universe grows. At first it gorges on the surrounding vacuum; later it
feeds more and more slowly; and after a span beyond all imagining---far longer
than $10^{100}$ years---it finally evaporates away. Is this our universe? No. And
the reason is the most useful clue in the whole picture.

This first universe was born straight from the void, with no parent outside it.
And a universe with no parent turns out to be a thin, plain place: it has roughly
the shortcomings of the textbook picture of cosmology. It makes a little more
matter than antimatter, but only a little. It has almost no \emph{dark matter} and
almost no \emph{dark energy}---the two great unseen ingredients that fill our own
sky. And measured across many such firstborn universes, they come out far larger
than ours. Something about being a firstborn leaves you sparse.

\section*{Our birth, and why our sky is dark}

But inside that first universe, exactly as inside ours, stars live and die and
black holes form---and each black hole does again what the void did once: it
bounces, and a new universe is born inside it. We are one of those. We are a child
universe, pressed into being through an enormous black hole inside our parent.

Picture the squeeze. The mass of an entire universe---very roughly $10^{53}$
kilograms---is forced through a throat no wider than a building, packed so tightly
that the whole of it would fit in a single room, and it happens in about a
thousandth of a billionth of a billionth of a second. At that density nothing
recognizable survives: not a torn teddy bear, not a rusty car, not an atom, not
even an atomic nucleus---all of it is melted into a featureless soup. What carries
across is not \emph{things} but a \emph{tally}: the faint surplus of matter over
antimatter, about one part in a billion, inherited from the parent and nudged a
little larger by the violence of the crossing.

And now the dark sky makes sense. Our parent's enormous mass sits just on the far
side of that birth-surface, and we feel its gravity \emph{before} its light can
ever reach us---a pull from matter we cannot see. That is dark matter. The parent
is still feeding matter through to us, slowly, and the gentle outward push of that
ongoing inflow is dark energy. Darkness is the harbinger of light. The parent's
own faint glow---the radiation any black hole gives off, as Stephen Hawking
showed---arrives smeared across our whole sky as the background hum that fills it.
The unevenness in all this inflow is the seed of everything: it gathers matter
into galaxies, into stars, into planets, and perhaps that is why we are here to
notice at all.

So we are not exotic. Count up all the universes and about nine in ten are
ordinary children like ours---half of them matter, half of them antimatter, each
with its own clock running its own way. We live at the most ordinary address in
the whole endless family.

\section*{Are we real, or a dream?}

There is one last doubt to face, and it is the sharpest. If universes are so
fantastically unlikely, isn't there a cheaper way to make a thinking being? Why
grow a whole cosmos when you could just let a mind flicker into existence by
accident, memories and all---or, what amounts to the same thing, run it as a
simulation on some machine that fell together by chance? Lay the three prices side
by side and the worry looks real:

\begin{center}
a stray mind: one in $10^{10^{69}}$\,;\qquad
a world-seed: one in $10^{10^{84}}$\,;\qquad
a ready-made big bang: one in $10^{10^{123}}$.
\end{center}

\noindent The stray mind is the cheapest of the three. So shouldn't we expect to
\emph{be} one---a lucky accident with invented memories, rather than real
inhabitants of a real world?

No---and seeing why is the keystone of the whole book. A mind is a thing that
\emph{thinks}, and thinking is work: it burns through order, it lays down
memories, it computes, and every bit of that needs a direction to time---a low
place to pour its waste and a high place to draw from. The void has no such
direction. It is perfectly still and perfectly even, with no arrow at all. So a
stray mind cannot form in the void: it would need the very thing---a flow of
time---that only a born world can provide. The cheap accident is not cheap from
the void. It is \emph{impossible} there.

The same blade cuts the ready-made big bang. A finished, perfectly arranged
universe handed over whole is also something that could only be \emph{placed},
never \emph{grown}---it assumes the order it was meant to explain. A picture of the
cosmos that can only shrug and say ``it simply happened to begin that way'' has,
without meaning to, declared us an accident with no maker: a simulation in all but
name. The way out is not to assume the order but to \emph{make} it---and the
bounce makes it, for free, every time, out of a cheap and generic seed.

So the choice was never ``cheap fake mind'' against ``expensive real world.'' Only
one of the two can come from nothing at all. The void cannot dream up a brain. It
can only give birth to a world---and then let the world, in its own long time,
grow the minds that wonder how they got here.

\bigskip
\noindent\emph{The rest of this book gives all of this its proper names, its
equations, and its tests---the bounce surface, the family tree of universes, the
inherited tally of matter, the background glow, the clocks that run at different
rates---and asks, at every step, how we could tell whether it is true. But the
story above is the whole of it.}
